%% mechanization -- what the Lean development checks, plain entailment as a
%% corollary, the mechanized bounds, where the formal route differs, what is
%% not mechanized, and the trusted base. Details and the correspondence with
%% the paper are in appendix-lean.

\section{Mechanization}
\label{sec:mechanization}

The artifact, a development in Lean~4 \citep{MouraUllrich21} with Mathlib
\citep{Mathlib20}, contains six decision procedures, one for each case of
\cref{thm:as-main}(1,2): existential entailment over $\BFT$ for every arity
$n\ge0$ and every variance, and over $\BF$ for every arity with $f$ covariant,
each over arbitrary, regular, and finite trees. Arity zero is included. The
procedures are computable definitions, not classical choices of
decidability instances, and each has a correctness theorem without premises:
none assumes satisfiability or positive arity. They take flat systems (\cref{sec:bg-constraints}) whose visible and hidden variables
are finite ordinals. Classical reasoning enters the proofs of the bounds and of
correctness, not the searches. Lean and Mathlib are both pinned to
\texttt{v4.35.0}, and the correctness theorems depend only on the axioms
\texttt{propext}, \texttt{Classical.choice}, and \texttt{Quot.sound}.

Plain entailment is the instance with no hidden variables and the consequent
$x\le y$ (\cref{def:bg-existential}). The artifact defines decision procedures
for plain entailment in the three tree classes with the same names and types
as those of the artifact of \citet{Zhao26}, and obtains them from the
existential procedures through \lean{nsse_iff_ensse} and its regular and
finite versions; it uses no other procedure for plain entailment. The same
reduction decides covariant plain entailment over $\BF$. The two artifacts
share, as a library, the definitions, the results on satisfiability, and the
hole search (\cref{sec:bg-readers}), but not the reduction of plain entailment to
holes, so the decidability of plain entailment is proved a second time, through
the characterization of \cref{sec:ch}.

Every search bound is an explicit function of the state counts of the
automata involved. Anchored compression (\cref{thm:mo-anchored}) is
mechanized with an elementary bound: \lean{Anchored.exists_short_posHole_elementary}
shortens a positive hole of a lettered monitor with a positive channel to
length at most \lean{Anchored.anchoredLengthBound}, and
\lean{Anchored.anchoredLengthBound_le} proves the estimate
$2^{2^{5\Xisharp(N,d)+6}}$, against $2^{2^{5\Xisharp(N,d)}}$ in
\cref{thm:mo-anchored}, with the same exponent $\Xisharp$. The dictionary bound
is $\sGsharp$ itself; the column count is larger than $\scolsharp$ because the
formal positive summary stores a second complete column summary where the
paper adds single fields (\cref{app:lean-diff}). The mortal bound evaluates
this function at the separator reader with all its states, not only the
reachable ones, and the non-mortal bound has the form of $\Lnm$ with a larger
pump capacity. Their maximum, \lean{Core.Words.Self.selfHoleBound}, bounds a
shortest hole of every instance with forced target heads, as in
\cref{thm:mo-decide}.

Three principal differences in the formal route are the following. For $\BF$
with $f$ covariant the formalization does not use the reduction of \cref{thm:tc-gadget}: a direct
covariant core represents a tree by its set of constructor positions
(\cref{lem:bg-bf-order}), proves three conditions in which preservation of
liveness takes the place of \CondU\ and no self events occur, reduces the
third to two-word holes without self guard, transfers to regular and finite
trees through regular Horn hulls and finite least solutions, and identifies
sets of constructor positions with $\Top$-free trees. For regular trees, (2) and (3) of
\cref{thm:tc-regular} are proved like \cref{thm:tc-finite}: the at most two
observations of \cref{lem:tc-two-seeds} become ground bounds; these bounds, and
in (2) the regular visible values, are expressed by literals over fresh
variables whose values are fixed by the graph equations of regular
presentations; and the extended system has a regular solution by (1).
\Cref{lem:sm-regular-closure} is not used. Condition \CondU\ is decided by the
hole search of \citet{Zhao26} on the unit reader with all its $N_{\mathrm U}$ states, without
the reductions of \cref{lem:un-reduced,lem:un-clock}.

Not mechanized are the flattening of type expressions
(\cref{lem:bg-flattening}); the bounds in terms of the number $s$ of variables,
that is, \cref{thm:as-main}(3,4), \cref{prop:as-bounds},
\cref{lem:mo-estimate,cor:mo-extracted}, and the closed form in
\cref{thm:un-decidable}; the effective countermodels of
\cref{lem:as-countermodel}, since the procedures return only the answer; the
results of \cref{sec:ap}; and the complexity results of \cref{sec:cx},
together with the polynomial-space and polynomial-time claims of
\cref{prop:sm-decide,lem:sm-mu,lem:sm-dead-states}, whose tests the artifact
performs by finite searches.

The trusted statement consists of the definitions of trees, the variance
order, finite terms and regular graphs with their interpretations as trees,
flat literals and their satisfaction, and the six entailment relations. Trees,
the order, and the two representations are those of \citet{Zhao26}, whose
characterization module relates the order to the greatest fixed point of the
constructor rule and shows that the representations denote the regular and
the finite trees. The kernel checks the proof terms; running the compiled
procedures also relies on the compiler. \Cref{app:lean} gives the public
statement, the module layout, the correspondence with the results of this
paper, and the differences in detail. The artifact is included in
\anon[the supplementary material]{\url{https://github.com/noti0na1/deci-ensse}}.
